% Generated by Listas/make_statements.py from pratica-lecture-08.md, item 22. Do not edit.
\textbf{Source.} Tadelis (2013), \emph{Game Theory: An Introduction}, Exercise 16.7 ``Intuitive Limit Pricing'', p.~338 {[}PDF p.~357{]}. Reproduced verbatim. Answer: none in the book.



\begin{quote}

\textbf{Intuitive Limit Pricing:} Show that of all the separating equilibria of the limit pricing--entry deterrence game in Section 16.2 only the best separating equilibrium does not fail the intuitive criterion.

\end{quote}



Note: the game of Tadelis §16.2 (pp.~323--332) is worked in \texttt{Notas/L08/L08-07-limit-pricing.md} (inverse demand \(5 - q\), costs 1 and 2, separating \(q_1^{1L} \in [2.62, 3.32]\)). Errata: Tadelis p.~331 prints the entrant's pooled-entry profit as \(\frac{7}{18}\), arithmetic gives \(\frac{7}{36}\); Def. 16.1 condition 2 reads \(v_2\) where \(v_1\) belongs.



\textbf{Game of Section 16.2 (pp.~323--324 {[}PDF pp.~342--343{]}), copied from the book because the exercise refers to it.}



\begin{quote}

Consider a market for some good that will last for two periods with demand \(P = 5 - q\) in each period. In period \(t = 1\) only one firm, the incumbent (player 1), is in the market, and this incumbent will continue to produce in period \(t = 2\). A potential entrant (player 2) observes what happened in period \(t = 1\) and then makes a choice of whether to enter the market at a fixed cost of \(\frac{1}{2}\) and compete against the incumbent in period \(t = 2\), or to remain out of the market. If player 2 stays out of the market then player 1 remains a monopolist in period 2, while if player 2 enters the market then they compete in a Cournot duopoly game. Each firm has marginal costs \(c_i\), where it is common knowledge that the entrant's marginal costs are \(c_2 = 2\). There is asymmetric information with respect to the costs (the type) of the incumbent. In particular the entrant knows only that \(c_1 \in \{1, 2\}\), so that the incumbent can be either more efficient than the entrant (have low {[}\(L\){]} costs of \(c_1 = 1\)) or as efficient (have high {[}\(H\){]} costs of \(c_1 = 2\)). It is common knowledge that \(\Pr\{c_1 = 1\} = \frac{1}{2}\), but only the incumbent knows his true costs.



The timing of this game, which captures the story, is described as follows:



\begin{enumerate}

\def\labelenumi{\arabic{enumi}.}

\tightlist

\item

  Nature chooses the type of player 1, \(c_1 \in \{1, 2\}\), each type with equal probability.

\item

  Player 1 observes his costs and decides how much to produce in the first period, \(q_1^1\), and the price in the first period is then \(P = 5 - q_1^1\).

\item

  Player 2 observes \(q_1^1\) and decides whether or not to enter at a fixed cost of \(F = \frac{1}{2}\).

\item

  If player 2 stays out then in period \(t = 2\) player 1 chooses \(q_1^2\) and the price is \(P = 5 - q_1^2\).

\item

  If player 2 enters then in period \(t = 2\) each player \(i\) simultaneously chooses \(q_i^2\) and the price is \(P = 5 - q_1^2 - q_2^2\) (Cournot competition).

\end{enumerate}

\end{quote}

